% ------------------------------------------------------------
% Page 79
% ------------------------------------------------------------

\begin{center}
    {\Large\bfseries La Société Cévenole de Carburants Forestiers}
\end{center}

\vspace{1.2cm}

L’histoire de la Société Cévenole des Carburants Forestiers, qui fut créée, près de
Meyrueis, pendant l’hiver 1942, est exemplaire. Il s’agissait, grâce à un brevet suisse,
de produire un charbon de bois adapté aux gazogènes. La construction des fours fut
entreprise au lieu-dit Alauze, à deux km de Meyrueis, dans la vallée de la Brèze. Le
directeur Georges Roland et le comptable Robert Cabané, ne tardèrent pas à recruter
de nombreux jeunes désireux de travailler et d’échapper au STO (Service du Travail
Obligatoire), comme le leur permettait le classement officiel de la Société Cévenole en
tant qu’« usine prioritaire ».

\vspace{0.45cm}

L’un des premiers à s’embaucher fut un réfractaire rescapé d’Aire de Côte. (maquis
installé à la Maison Forestière d’Aire de Côte, non loin de l’Aigoual, décimé par une
attaque de chasseurs parachutistes allemands le 1\textsuperscript{er} Juillet 1943).
Puis ce furent sept ou huit polonais placés en résidence surveillée à l’Hôtel de France.
Des membres de la famille du directeur, un prisonnier évadé, des requis nîmois du STO,
deux Alsaciens, vinrent les rejoindre, tous bénéficiant des cartes de ravitaillement
obtenues par des contacts personnels auprès des services gardois compétents.

\vspace{0.45cm}

Dans certains cas, l’inscription sur les registres d’embauche était fictive et servait de
couverture aux bénéficiaires. Dans d’autres cas, au contraire, des réfractaires trop
exposés faisaient le travail sans être inscrits sur les registres. « Quelle gymnastique
de chiffres ! » nous dit le comptable chargé de faire viser le livre de paye à Mende.
C’est ainsi que, sur les trente cinq employés réguliers de la Société Cévenole,
presque une vingtaine peuvent être considérés comme des « réfugiés » plus ou moins
clandestins.

\vspace{0.45cm}

Mais cette « planque » connut aussi plusieurs drames. Dans la nuit du 29 au 30 mars
1944, des camions allemands vinrent enlever les Polonais de l’Hôtel de France et
l’on n’entendit plus parler d’eux.

\vspace{0.45cm}

Parmi les employés, le jeune Noguès, qui avait cessé le travail pour suivre le maquis
Bir Hakeim, fut capturé à la Parade et fusillé. Un autre employé de la Société
cévenole, le jeune Fages, qui s’était aussi engagé à Bir Hakeim, fut grièvement
blessé au combat, recueilli dans la vallée de la Jonte, où il avait pu se traîner,
transporté dans la voiture de liaison de l’usine et soigné, d’abord au Moulin d’Ayres
chez les Loupiac, puis à Vébron, chez le pasteur Chazel.

\vfill

\begin{flushright}
Cévennes terre de refuge 1940-1944\\
Presses du Languedoc P. 290-291
\end{flushright}

\clearpage
